\section{引言}
人体活动识别旨在根据传感器观测判断人的运动状态，是移动健康与人机交互中的基础任务。智能手机惯性传感器能够连续记录加速度和角速度，但同一活动在不同人的动作幅度、节奏与姿态下可能产生不同信号。因此，评价模型时需要关注其对未参与训练的受试者的识别能力，不能仅依赖同一批受试者内部的随机窗口划分。

UCI HAR 为此提供了同时包含人工特征和惯性信号窗口的公开数据资源\cite{uci,anguita2013}。人工特征将时域与频域统计信息压缩为固定维数表示，适合传统分类器；时序网络则能够直接从窗口序列中学习局部变化模式。支持向量机与随机森林分别代表核方法和树集成方法\cite{cortes1995,breiman2001}。可穿戴活动识别中的卷积与循环网络研究说明了从信号学习表示的可行性\cite{ordonez2016}，但不同数据集、输入通道及划分协议下的分数不能直接作为统一排名依据。

本文考察人工特征表示与学习得到的时序表示是否具有互补性。以逻辑回归为强基准，构建常见的 MLP--CNN 双分支拼接网络，并在同一组外层受试者划分上比较六类模型。研究包括人工特征降维对照、融合模型的随机种子复验，以及冻结模型后的官方独立测试。重点回答三个问题：融合是否改善验证和测试表现；降维是否带来性能与计算开销的折中；总体指标的变化是否同时惠及所有类别和受试者。本文采用的拼接机制属于通用表示融合方法，研究价值主要体现在可追溯的实验流程和对收益边界的分析。
